% !TeX root = ../../Book2.tex
% 纲要标题：## 第13章　半单模与 Artin–Wedderburn 理论
% 纲要位置：计划纲要.md，第 1462 行。

\chapter{半单模与 Artin–Wedderburn 理论}
\label{b2:ch:13}
\BookChapterNumber{b2:ch:13}

\begin{chapterabstract}
半单模可以分解为简单模的直和，等价地，它的每个子模都有补模。本章先证明这一判据及子商稳定性，再由正则模的有限简单分解建立 Artin–Wedderburn 定理。证明通过简单模之间的同态得到矩阵环，并逐步核对左模约定中的反环以及分解的唯一性。最后讨论有限维代数和群代数的基本例子，以乘法群的类方程和分圆多项式证明有限除环的交换性。
\end{chapterabstract}

\begin{learninggoals}
\begin{itemize}
\item 证明简单模直和、简单子模之和及所有子模有补三个条件等价，并处理无限直和的选择步骤。
\item 用 Schur 引理计算简单类型之间的 Hom，区分除环结论与代数闭域上的标量结论。
\item 从正则模推导矩阵环乘积分解，证明参数唯一性，并正确处理左右模与反环。
\item 比较半单性与半本原性，计算有限维和有限半单环的简单模及维数，证明 Wedderburn 小定理。
\end{itemize}
\end{learninggoals}

向量空间的每个子空间都有补空间，但 \(\mathbb Z/4\mathbb Z\) 中的子群 \(2\mathbb Z/4\mathbb Z\) 没有补子群：任何阶为二的子群仍是它自身。它的合成列有两个 \(\mathbb Z/2\mathbb Z\) 商因子，却不能把这两个因子拆成直和。\cref{b2:thm:module-jordan-holder}确定的是逐层出现的简单商及其重数；本章的半单性则要求简单模已作为原模内部的直和分量存在。二者的区别在于连接这些商层的扩张是否分裂。

\ProseParagraphBreak

从这个补模条件出发，可以先在模上辨认半单性，再问一个环何时使自身作为模半单。对于除环 \(D\)，矩阵环 \(M_n(D)\) 的左正则模可按列分解为 \(n\) 个同构的简单列模。一般半单环的分类要从正则模的简单类型及重数恢复这些矩阵块，并弄清每块所用的除环；这需要计算自同态环，不能只比较维数。

\ProseParagraphBreak

本章使用模的商、直和、正合列及自同态环；有限长度的术语沿用第十二章。无限直和的补模构造使用第一卷定理 A.1.4 的 Zorn 引理。Wedderburn 小定理还用到第一卷的有限群类方程和分圆多项式的整数因式分解。

\ProseParagraphBreak
有限维表示中的 Schur 引理和半单分解可参阅 Etingof 等人的《Introduction to Representation Theory》\cite[Proposition 2.3.9 and Corollary 2.3.10, pp. 11--12; Proposition 3.1.4 and Remark 3.1.5, pp. 42--43]{EtingofEtAl2011RepresentationTheory}；该书\cite[Proposition 3.5.8, p. 48]{EtingofEtAl2011RepresentationTheory}给出代数闭域上有限维代数的对应结构定理。本章先处理一般环与任意大小的半单模，再在有限维代数的应用中加上基域假设。

% !TeX root = ../../Book2.tex
% 纲要标题：### 13.1 简单模与半单模
% 纲要位置：计划纲要.md，第 1464 行。

\section{简单模与半单模}
\label{b2:sec:1301}

一个模可以有很短的合成列，却仍不能把各个商因子直接拆开。半单性要求更强：组成模本身的简单部分已经作为直和存在，所有子模都能找到补模。本节允许无限直和，不预设模有限生成。环仍含单位元，模仍为幺性左模；简单模总要求非零。

% 纲要要点（供目录核对）：
%
% * 简单对象和简单模的直和
% * 半单模的子模与商模
% * Schur 引理及除环；逆命题与有限维假设的反例
%

\subsection{简单对象和简单模的直和}
\label{b2:subsec:1301:01}

\begin{definition}\label{b2:def:simple-module}
设 \(R\) 为含幺环，所论左 \(R\) 模均为幺性模。非零左 \(R\) 模 \(S\) 只有 \(0,S\) 两个子模时，称 \(S\) 为简单模\termidx[jiandan-mo]{简单模}。左 \(R\) 模 \(M\) 可写为一族简单子模的内部直和
\[
M=\bigoplus_{i\in I}S_i,
\]
时，称 \(M\) 为\term[bandan-mo]{半单模}[semisimple module]。
\label{b2:def:semisimple-module}
允许 \(I=\varnothing\)，因此零模半单，但不是简单模。
\end{definition}

例如，域 \(k\) 上的简单模是一维向量空间，每个向量空间选基后都是一维空间的直和，所以半单。对素数 \(p\)，整数模 \(\mathbb Z/p\mathbb Z\) 简单，而短正合列
\[
0\longrightarrow\mathbb Z/p\mathbb Z\xrightarrow{\bar a\mapsto\overline{pa}}
\mathbb Z/p^2\mathbb Z\longrightarrow\mathbb Z/p\mathbb Z\longrightarrow0
\]
不分裂。右端被 \(p\) 零化，所以从右端到中项的任何同态，其像都落在中项被 \(p\) 零化的子模 \(p\mathbb Z/p^2\mathbb Z\) 中；商映射在这个子模上为零，故该同态不可能是截面。这个唯一的非零真子模没有补模；它也是中项唯一的简单子模，所以中项不能写成简单子模的直和，\(\mathbb Z/p^2\mathbb Z\) 不是半单模。它的两个简单合成因子是在商层中出现的，并非模内已经存在的两个简单直和分量。

\ProseParagraphBreak
模范畴中的简单对象就是上述简单模。给定 \(0\ne s\in S\)，子模 \(Rs\) 非零，故 \(Rs=S\)。因此每个简单模都循环，且形如 \(R/L\)，其中 \(L\) 是极大左理想：满射 \(R\to S\)、\(r\mapsto rs\) 的核为 \(L\)，中间子模对应迫使 \(L\) 极大。反过来，极大左理想的商模只有零子模和全模，因而简单。这里 \(L\) 只要求左理想，不必是双边理想，商对象通常只是模。

\ProseParagraphBreak
例如，令 \(k\) 为域、\(R=M_2(k)\)，并取第一列为零的矩阵组成的集合 \(L\)。左乘任意矩阵仍使第一列为零，所以 \(L\) 是左理想。映射 \(R\to k^2\) 取矩阵的第一列，是满射左 \(R\) 模同态，核恰为 \(L\)。列模 \(k^2\) 简单：给定非零列向量 \(v\) 与任意 \(w\)，选 \(v_j\ne0\)，令矩阵的第 \(j\) 列为 \(w/v_j\)、另一列为零，就有 \(Av=w\)。因此 \(R/L\cong k^2\) 简单，\(L\) 是极大左理想。但是，若 \(E_{ij}\) 表示第 \((i,j)\) 个位置为 \(1\)、其余位置为零的矩阵，则 \(E_{12}\in L\)，而 \(E_{12}E_{21}=E_{11}\notin L\)；故 \(L\) 不是双边理想，\(R/L\) 不能以代表元乘法成为商环。

\ProseParagraphBreak
“简单子模之和”暂时不要求和为直和。以下定理说明，只要简单子模的和已经覆盖整个模，总可以从中抽出一个直和分解；同时，它给出不依赖所选简单分解的半单性判据。

\begin{theorem}[半单模的等价刻画]\label{b2:thm:semisimple-module-characterizations}
设 \(R\) 为含幺环，\(M\) 为幺性左 \(R\) 模。下列条件等价：
\begin{equivalences}
\item \(M\) 是简单子模的直和。
\item \(M\) 是其简单子模之和。
\item 对每个子模 \(N\subseteq M\)，存在子模 \(K\) 使 \(M=N\oplus K\)。
\end{equivalences}
\end{theorem}
\begin{proof}
第一条立即蕴含第二条。设 \(M=\sum_{i\in I}S_i\)，其中每个 \(S_i\) 简单，固定 \(N\subseteq M\)。为了给 \(N\) 找补模，考察所有子集 \(J\subseteq I\)，使 \(\sum_{j\in J}S_j\) 是直和且与 \(N\) 交为零，按包含排序。空集满足条件。链的并仍满足这两个条件。确实，任何有限个指标都同属链中的某一项，所以这些分量之间若有和为零的关系，直和性使每个分量为零；这些分量之和若落入 \(N\)，这一项与 \(N\) 零交又使该和为零。由 Zorn 引理，取极大的 \(J\)，令 \(K=\bigoplus_{j\in J}S_j\)。

若 \(N+K\ne M\)，某个 \(S_i\) 不包含于 \(N+K\)。交 \(S_i\cap(N+K)\) 是简单模 \(S_i\) 的子模，只能为零或整个 \(S_i\)；后一情形已被排除，故此交为零。于是 \(K+S_i\) 仍为直和；若 \(k+s\in N\)，其中 \(k\in K,s\in S_i\)，则 \(s\in S_i\cap(N+K)=0\)，继而 \(k\in K\cap N=0\)。所以可将 \(i\) 加入 \(J\)，保持两个条件，与极大性矛盾。因此 \(M=N\oplus K\)，得到第三条。特别地，取 \(N=0\) 便已经得到第一条。

最后假设第三条。要把简单商层变成模内的简单子模，先注意补模条件传给每个子模 \(U\)：若 \(L\subseteq U\) 且 \(M=L\oplus C\)，将 \(u\in U\) 写成 \(u=l+c\) 后，由 \(l\in L\subseteq U\) 得 \(c\in U\cap C\)，所以 \(U=L\oplus(U\cap C)\)。取非零子模 \(U\) 中的 \(0\ne x\)。循环模 \(Rx\) 的真子模按包含排序；任一链的并仍不含生成元 \(x\)，故仍为真子模。Zorn 引理给出极大真子模 \(L\)，于是 \(Rx/L\) 简单。由传给 \(Rx\) 的补模条件，\(Rx=L\oplus S\)，商映射限制为同构 \(S\to Rx/L\)。这样，原先只在商模中出现的简单模被实现为 \(U\) 内的非零简单子模 \(S\)。

令 \(N\) 为 \(M\) 的全部简单子模之和。如果 \(N\ne M\)，取非零补模 \(K\)。刚才的结论给出 \(K\) 的简单子模 \(S\)，但 \(S\subseteq N\) 又使 \(S\subseteq N\cap K=0\)，矛盾。因此 \(N=M\)，得到第二条，完成等价性证明。
\end{proof}

\subsection{半单模的子模与商模}
\label{b2:subsec:1301:02}

\begin{proposition}\label{b2:prop:semisimple-subquotients}
设 \(R\) 为含幺环，所论左 \(R\) 模均为幺性模。半单左 \(R\) 模的每个子模与商模都半单，任意一族半单左 \(R\) 模的直和也半单。半单左 \(R\) 模有限生成，当且仅当它是有限多个简单子模的直和。对半单左 \(R\) 模 \(M\) 的子模 \(N\)，短正合列
\[
0\longrightarrow N\longrightarrow M\longrightarrow M/N\longrightarrow0
\]
分裂。
\end{proposition}
\begin{proof}
由\cref{b2:thm:semisimple-module-characterizations}，写 \(M=N\oplus K\)。取一个简单分解 \(M=\bigoplus_iS_i\)，投影 \(p:M\to N\) 将每个 \(S_i\) 映为零或简单模：限制映射的核为 \(0\) 或 \(S_i\)，非零像因此同构于 \(S_i\)。又 \(N=p(M)=\sum_i p(S_i)\)，由等价刻画得 \(N\) 半单。这里得到的是投影后的简单子模，不必是原来 \(S_i\) 中某些分量的和。同理投影到 \(K\) 说明 \(K\) 半单，而商映射限制为同构 \(K\to M/N\)，故商模半单，其逆后接 \(K\hookrightarrow M\) 就给出商映射的截面。

若各 \(M_a\) 都已分解为简单模直和，把全部指标并在一起，仍是直接和：任一元素只有有限多个非零外层分量，每个分量又只有有限多个非零内层分量。因此所有简单分量合成所需直和分解。

若 \(m_1,\ldots,m_t\) 生成半单模 \(M=\bigoplus_{i\in I}S_i\)，它们的支集之并是有限集 \(F\subseteq I\)。每个 \(S_i\) 是子模，故对生成元作标量乘法及求和不会产生 \(F\) 之外的直和坐标。因此 \(M=\bigoplus_{i\in F}S_i\)，只含有限多个简单分量。反过来，每个简单模由任一非零元素循环生成，有限个简单模的直和便由这些元素有限生成；零模对应空直和。
\end{proof}

两端半单不保证中项半单：前面的 \(\mathbb Z/p^2\mathbb Z\) 正是两个简单模之间不分裂的扩张。合成因子的名单没有记录这类扩张是否拆开。另一方面，命题中的有限生成判据将把半单环的简单分解限制为有限直和。

\subsection{Schur 引理及除环}
\label{b2:subsec:1301:03}

\begin{lemma}[Schur]\label{b2:lem:schur}
设 \(R\) 为含幺环，\(S,T\) 为幺性简单左 \(R\) 模。任何非零同态 \(f:S\to T\) 都是同构。因此若 \(S\not\cong T\)，则 \(\operatorname{Hom}_R(S,T)=0\)；而
\[
D=\operatorname{End}_R(S)
\]
是除环，其乘法为映射复合。
\end{lemma}
\begin{proof}
\(\ker f\) 是 \(S\) 的子模；因 \(f\ne0\)，它不能为 \(S\)，所以为零。\(\operatorname{im}f\) 是 \(T\) 的非零子模，所以为 \(T\)。故 \(f\) 双射；模同态的逆仍是模同态。取 \(S=T\)，每个非零自同态都可逆，恒等映射因 \(S\ne0\) 而非零，所以自同态环是除环。
\end{proof}

这一结果称为\term[Schur-yinli]{Schur 引理}[Schur's lemma]。结论是除环，不能在一般环上把它直接写成标量域。例如除环 \(D\) 的左正则模简单，而\cref{b2:prop:regular-hom-endomorphism}给出其自同态环为 \(D^{\mathrm{op}}\)。对含幺 \(k\) 代数 \(A\) 的非零幺性模 \(S\)，标量映射给出 \(k\) 到 \(\operatorname{End}_A(S)\) 中心的嵌入，但这个嵌入未必满射。例如 \(\mathbb C\) 作为实代数作用在自身上，也是简单模，其自同态环为 \(\mathbb C\)，大于原标量域 \(\mathbb R\)。

\begin{corollary}\label{b2:cor:schur-algebraically-closed}
设 \(k\) 为代数闭域、\(A\) 为结合含幺 \(k\) 代数，\(S\) 为非零、有限 \(k\) 维的幺性简单左 \(A\) 模，则
\(\operatorname{End}_A(S)=k\operatorname{id}_S\)。
\end{corollary}
\begin{proof}
取 \(f\in\operatorname{End}_A(S)\)。由于 \(S\) 有限维且 \(k\) 代数闭，\(f\) 的特征多项式有根 \(\lambda\in k\)，所以 \(f-\lambda\operatorname{id}_S\) 有非零核。该映射仍 \(A\) 线性，因为 \(k\) 的作用在中心。由\cref{b2:lem:schur}，它不可能是非零同态，故 \(f=\lambda\operatorname{id}_S\)。反过来，每个这样的标量映射当然与 \(A\) 作用相容。
\end{proof}

例如，若 \(S,T\) 非同构且简单，自同态 \(f:S\oplus T\to S\oplus T\) 的四个分块分别是两个对角自同态及 \(S\to T\)、\(T\to S\) 的映射。后两者由 Schur 引理为零，所以分块形式为
\[
\begin{pmatrix}a&0\\0&b\end{pmatrix},\qquad
a\in\operatorname{End}_R(S),\quad b\in\operatorname{End}_R(T),
\]
自同态环因此是 \(\operatorname{End}_R(S)\times\operatorname{End}_R(T)\)。如果改为 \(S\oplus S\)，四个分块都属于 \(D=\operatorname{End}_R(S)\)，映射可以写为
\[
f(s_1,s_2)=\bigl(a_{11}(s_1)+a_{12}(s_2),\,
a_{21}(s_1)+a_{22}(s_2)\bigr),\qquad a_{ij}\in D.
\]
复合时先施加后一映射，再施加前一映射，四个条目按矩阵乘法组合，故 \(\operatorname{End}_R(S\oplus S)\cong M_2(D)\)。这一区别正是 Artin–Wedderburn 分解中“不同块”与“同一块内的矩阵”各自出现的原因。

\begin{proposition}\label{b2:prop:schur-hypothesis-counterexamples}
幺性整数模 \(\mathbb Q\) 不是简单模，但 \(\operatorname{End}_{\mathbb Z}(\mathbb Q)\cong\mathbb Q\) 是除环。另设 \(k\) 为代数闭域、\(t\) 为 \(k\) 上的不定元，\(A=k(t)\)。左正则模 \(S={}_AA\) 是幺性简单左 \(A\) 模，具有无限 \(k\) 维数，且
\[
\operatorname{End}_A(S)\cong k(t)\supsetneq k.
\]
因此自同态环为除环不蕴含模简单，而代数闭基域上的简单模若没有有限维假设，其自同态也未必全是基域标量。
\end{proposition}
\begin{proof}
若 \(f:\mathbb Q\to\mathbb Q\) 为 \(\mathbb Z\) 线性，令 \(q=f(1)\)。对 \(m\in\mathbb Z\) 和正整数 \(n\)，有 \(nf(m/n)=f(m)=mq\)，故 \(f(m/n)=mq/n\)。因此 \(f\) 恰为乘以 \(q\) 的映射。反过来，乘以任意 \(q\in\mathbb Q\) 的映射都 \(\mathbb Z\) 线性，加法与复合分别对应有理数的加法与乘法，于是得到环同构 \(\operatorname{End}_{\mathbb Z}(\mathbb Q)\cong\mathbb Q\)。但 \(\mathbb Z\) 是 \(\mathbb Q\) 的非零真子模，所以 \(\mathbb Q\) 不简单。

因为 \(A=k(t)\) 是域，任意非零左理想含有一个可逆元，从而为全环，故其左正则模简单。向量 \(1,t,t^2,\ldots\) 在 \(k\) 上线性无关，所以 \(S\) 无限维。任意 \(A\) 线性自同态满足 \(f(a)=af(1)\)，而任意乘以 \(c\in A\) 的映射都 \(A\) 线性；由于 \(A\) 交换，这给出 \(\operatorname{End}_A(S)\cong A\)。其中乘以 \(t\) 的映射不是 \(k\) 标量映射：若它等于 \(\lambda\operatorname{id}_S\)，在 \(1\) 上求值就会得到 \(t=\lambda\in k\)，矛盾。事实上该映射没有 \(k\) 中的特征值，因为 \((t-\lambda)a=0\) 对每个 \(\lambda\in k\) 都迫使 \(a=0\)。这正说明前一推论中有限维假设保证特征值存在的作用。
\end{proof}

% !TeX root = ../../Book2.tex
% 纲要标题：### 13.2 半单环
% 纲要位置：计划纲要.md，第 1470 行。

\section{半单环}
\label{b2:sec:1302}

半单模可以定义在任意环上；半单环则要求环自身的正则模已经半单。因为每个模都由自由模取商得到，这个看似只涉及一个模的要求，实际上会迫使整个模范畴中的子模都能分裂。它还蕴含强烈的有限性，不能与仅仅“没有非零元素同时零化所有简单模”混淆。

% 纲要要点（供目录核对）：
%
% * 正则模半单的等价条件
% * 所有模的投射性与内射性
% * 半单性不等于半本原性
%

\subsection{正则模半单的等价条件}
\label{b2:subsec:1302:01}

\begin{definition}\label{b2:def:semisimple-ring}
设 \(R\) 为含幺环。若左正则模 \({}_RR\) 是半单模，就称 \(R\) 为\term[bandan-huan]{半单环}[semisimple ring]。零环的正则模为零，按此约定也是半单环。\secref{b2:sec:1303}将证明该定义与右正则模半单等价。
\end{definition}

\begin{theorem}\label{b2:thm:semisimple-ring-characterizations}
设 \(R\) 为含幺环，所论左 \(R\) 模均为幺性模。下列条件等价：
\begin{equivalences}
\item 左正则模 \({}_RR\) 半单。
\item 每个左理想都是 \({}_RR\) 的直和因子。
\item 每个左 \(R\) 模都是半单模。
\item 每条左 \(R\) 模的短正合列都分裂。
\end{equivalences}
若这些条件成立，\({}_RR\) 是有限多个简单左理想的直和，因而具有有限长度，特别是左 Noether 且左 Artin。
\end{theorem}
\begin{proof}
第一、二条的等价是\cref{b2:thm:semisimple-module-characterizations}应用于 \({}_RR\)，因为正则模的子模恰是左理想。假设第一条，要从这个正则模推广到任意左模 \(M\)，可先把它的生成元放进自由模：自由模 \(R^{(I)}\) 是若干份 \({}_RR\) 的直和，由\cref{b2:prop:semisimple-subquotients}仍半单。任意左模 \(M\) 都有满射 \(R^{(I)}\to M\)，例如按 \(M\) 的元素取自由生成元并将其送到对应元素。再用同一命题中商模保持半单的结论，得到第三条。第三条当然包含第一条。

若第三条成立，短正合列 \(0\to A\xrightarrow{i}B\xrightarrow{q}C\to0\) 中，\(i(A)\) 在 \(B\) 内有补模 \(K\)。\(q|_K\) 单射是因为 \(K\cap\ker q=0\)，满射是因为每个 \(b\) 可写成 \(i(a)+k\)。其逆后接 \(K\hookrightarrow B\) 便是截面，所以序列分裂。反过来，将第四条用于每条 \(0\to N\to M\to M/N\to0\)，由\cref{b2:thm:split-exact-sequence}得到每个子模都有补模，故每个 \(M\) 半单。

最后，\(R\) 由 \(1\) 生成，\cref{b2:prop:semisimple-subquotients}中的有限生成判据便使简单分量只有有限多个。具体地，若 \(R=\bigoplus_{i\in I}S_i\)，则 \(1=\sum_{i\in F}s_i\) 只涉及有限集 \(F\subseteq I\)。每个 \(S_i\) 是左理想，所以对任意 \(r\in R\)，有 \(r=r1=\sum_{i\in F}rs_i\in\bigoplus_{i\in F}S_i\)。因此这有限个分量已覆盖全部 \(R\)。逐个加上它们得到有限合成列，故 \({}_RR\) 有限长度。由\cref{b2:thm:finite-length-chain-conditions}，左理想的升链与降链都稳定，得到两项有限性结论。
\end{proof}

在此意义下，“半单 Artin 环”中的左 Artin 条件已由半单性推出。对域 \(k\)，无限直积环 \(R=\prod_{i\ge1}k\) 不会仅因每个因子是域就自动成为半单环：其正则模含有严格下降的左理想链
\[
R\supsetneq\{x:x_1=0\}\supsetneq
\{x:x_1=x_2=0\}\supsetneq\cdots.
\]
每一步都可由仅在相应坐标为 \(1\) 的元素见出严格包含。定理保证半单环必左 Artin，所以这个环不是半单环。它的各个坐标理想虽是简单模，但它们的直和只包含有限支集的元素，不能覆盖整个正则模。

\ProseParagraphBreak
这里讨论的是无限直积环在自身乘法下的正则模，不能据此断言同一个加法群在别的标量作用下也不半单。例如 \(\prod_{i\ge1}\mathbb F_2\) 作为整数模被 \(2\) 零化，因而是一个 \(\mathbb F_2\) 向量空间。选取向量空间基后，它就是一维 \(\mathbb F_2\) 空间的直和，每个分量作为整数模都同构于简单模 \(\mathbb Z/2\mathbb Z\)，故这个整数模半单。这里选取的基不是各个坐标向量组成的集合；后者只生成有限支集的向量。半单性取决于底环及其作用。

\subsection{所有模的投射性与内射性}
\label{b2:subsec:1302:02}

\begin{proposition}\label{b2:prop:semisimple-projective-injective}
设 \(R\) 为含幺环，所论左 \(R\) 模均为幺性模。若 \(R\) 半单，则每个左 \(R\) 模既投射又内射。反过来，若每个左 \(R\) 模都投射，或每个左 \(R\) 模都内射，则 \(R\) 半单。
\end{proposition}
\begin{proof}
设 \(R\) 半单。给定满射 \(q:B\to C\)，其核所确定的短正合列分裂，故有截面 \(s:C\to B\)。对任意 \(f:P\to C\)，映射 \(sf:P\to B\) 是提升，说明任意 \(P\) 投射。给定单射 \(i:A\to B\)，分裂性给出缩回 \(r:B\to A\)；任意 \(f:A\to E\) 都由 \(fr:B\to E\) 延拓，说明任意 \(E\) 内射。

反过来，若所有模都投射，对每条短正合列的商模使用提升性质，将其恒等映射提升为截面，所以所有短正合列都分裂。若所有模都内射，则将子模的恒等映射沿包含延拓，得到缩回，也使每条短正合列分裂。两种情形均由\cref{b2:thm:semisimple-ring-characterizations}推出半单性。
\end{proof}

这里要求“每个左模”都具有性质。给定模 \(M\) 的半单性只保证 \(M\) 自身的子模有补模；投射性却要求对任意满射提升从 \(M\) 出发的同态，内射性要求沿任意单射延拓以 \(M\) 为目标的同态，条件涉及其他模。因而一个给定模半单，不意味着它在原环上投射或内射。例如对素数 \(p\)，简单整数模 \(\mathbb Z/p\mathbb Z\) 不是投射模，因为 \(\mathbb Z\to\mathbb Z/p\mathbb Z\) 没有截面：它的任何截面都会把非零有限阶元素嵌入无挠模 \(\mathbb Z\)。它也不是内射模，因为从 \(p\mathbb Z\) 到它的同态 \(p\mapsto\bar1\) 不能延拓到 \(\mathbb Z\)：任何延拓都会把 \(p\) 送到 \(p\) 倍的某个元素，即零。

\ProseParagraphBreak
半单环上的直和分解通常不唯一到具体子模。例如 \(k^2\) 中 \(ke_1\) 的补空间可取任意 \(k(e_2+ae_1)\)。半单性保证补模存在，结构定理则描述简单类型及其重数；这两者都不要求为每个子模指定唯一补模。

\subsection{半单性与半本原性}
\label{b2:subsec:1302:03}

每个简单模 \(S\) 给出一个双边理想 \(\operatorname{Ann}_R(S)\)，由同时零化 \(S\) 全部元素的标量组成。把这些理想取交，可以检测一个环的所有简单模共同看不见哪些元素。

\ProseParagraphBreak
必须区别零化整个 \(S\) 与零化一个选定向量。对 \(0\ne s\in S\)，满射 \(R\to S\)、\(r\mapsto rs\) 的核 \(L_s=\{\,r\in R\mid rs=0\,\}\) 是极大左理想，而 \(\operatorname{Ann}_R(S)=\bigcap_{s\ne0}L_s\) 要求所有向量同时被零化。即使 \(S\cong R/L_s\)，也不能一般地写成 \(\operatorname{Ann}_R(S)=L_s\)。前面的 \(M_2(k)\) 列模就是例子：对 \(s=(1,0)^{\mathsf t}\)，\(L_s\) 是第一列为零的矩阵组成的非零左理想；但一个矩阵若零化全部列向量，其每一列都为零，故 \(\operatorname{Ann}_{M_2(k)}(k^2)=0\)。

\ProseParagraphBreak
所有简单模的同构类可以用一个集合表示：每个简单模都同构于某个 \(R/L\)，而极大左理想组成 \(R\) 的幂集中的一个集合。从这些商模中每个同构类选一个代表，便得到所需代表族；同构的模具有相同的零化子。

\begin{proposition}\label{b2:prop:simple-annihilator-intersection}
设 \(R\) 为非零含幺环，\(\Sigma\) 为幺性简单左 \(R\) 模的一组同构类代表。有
\[
\bigcap_{S\in\Sigma}\operatorname{Ann}_R(S)
=\bigcap_{L\text{ 为 }R\text{ 的极大左理想}}L.
\]
\end{proposition}
\begin{proof}
若 \(a\) 零化每个简单模，对每个极大左理想 \(L\)，商模 \(R/L\) 简单，故 \(a(1+L)=0\)，即 \(a\in L\)。这证明左边包含于右边。

反过来，设 \(a\) 属于全部极大左理想。固定任意简单模 \(S\)，但让 \(0\ne s\in S\) 遍历其所有非零向量。对每个这样的 \(s\)，满射 \(R\to S\)、\(r\mapsto rs\) 的核 \(L_s\) 都是极大左理想，因而包含 \(a\)，所以 \(as=0\)。这里可能需要随 \(s\) 改换 \(L_s\)，但 \(a\) 属于全部极大左理想的假设对每个 \(s\) 都适用。零向量也被零化，于是 \(a\in\operatorname{Ann}_R(S)\)，得到反向包含。
\end{proof}

若含幺环 \(R\) 的全部幺性简单左模的零化子之交为零，则称 \(R\) 具有\term[banbenyuan-xing]{半本原性}[semiprimitivity]。这就是说，简单模的作用族共同忠实：任何非零 \(a\in R\) 都在某个简单模的某个向量上作用非零；单个简单模的作用则未必忠实。第十四章将把这个交作为 Jacobson 根研究。

\ProseParagraphBreak
半单环一定满足这个条件。若 \(a\) 零化所有简单模，就零化左正则模的每个简单直和分量，因而零化整个 \(R\)；特别是 \(a=a1=0\)。但反过来不成立：\(\mathbb Z\) 的简单模为全部 \(\mathbb Z/p\mathbb Z\)，零化子为 \(p\mathbb Z\)，且
\[
\bigcap_{p\text{ 素数}}p\mathbb Z=0.
\]
若非零整数 \(a\) 属于此交，取 \(|a|+1\) 的任一素因子 \(p\)，它不整除 \(a\)，立即矛盾。但短正合列
\[
0\longrightarrow2\mathbb Z\longrightarrow\mathbb Z
\longrightarrow\mathbb Z/2\mathbb Z\longrightarrow0
\]
不分裂：若 \(2\mathbb Z\) 有补模 \(K\)，则 \(K\cong\mathbb Z/2\mathbb Z\)，而 \(\mathbb Z\) 中没有非零有限阶元素，不可能有这样的 \(K\)。因此 \(\mathbb Z\) 不半单。半本原性保证非零环元素能被某个简单作用检测，却没有保证子模与商模之间的扩张分裂。

\ProseParagraphBreak
零环没有非零简单模。若把空族理想的交按约定取为全环，上述交在零环中仍是零，半单性与半本原性在这个退化情形中均成立。后续结构定理将把零环对应于没有矩阵因子的空乘积，避免把它误列为除环。

% !TeX root = ../../Book2.tex
% 纲要标题：### 13.3 Artin–Wedderburn 定理
% 纲要位置：计划纲要.md，第 1476 行。

\section{Artin–Wedderburn 定理}
\label{b2:sec:1303}

半单环的正则模已分成有限多个简单部分。接下来考虑这些部分之间的全部同态：不同简单类型之间没有非零同态，同一类型的多个副本之间则组成矩阵。由于左正则模的自同态由右乘给出，最后必须通过反环把这个矩阵描述转回原环。

% 纲要要点（供目录核对）：
%
% * 半单环的矩阵块分解
% * 单 Artin 环与除环
% * 分解的唯一性、中心幂等元及左右约定
%

\subsection{半单环的矩阵块分解}
\label{b2:subsec:1303:01}

\begin{theorem}[Artin–Wedderburn]\label{b2:thm:artin-wedderburn}
设 \(R\) 为含幺环。\(R\) 半单，当且仅当存在有限多个除环 \(D_1,\ldots,D_t\) 与正整数 \(n_1,\ldots,n_t\)，使
\begin{equation}\label{b2:eq:artin-wedderburn}
R\cong\prod_{i=1}^tM_{n_i}(D_i).
\end{equation}
零环对应 \(t=0\) 的空乘积。每个这样的环也右半单，且左右均为 Artin 环和 Noether 环。
\end{theorem}
\begin{proof}
设 \(R\ne0\) 半单。由\cref{b2:thm:semisimple-ring-characterizations}，将正则模的有限简单分解按同构类型合并，写成
\[
{}_RR\cong\bigoplus_{i=1}^t S_i^{n_i},
\]
其中 \(S_i\) 两两非同构，\(n_i>0\)。令 \(E_i=\operatorname{End}_R(S_i)\)，由\cref{b2:lem:schur}它是除环。不同类型间的 Hom 为零：从 \(S_j^{n_j}\) 到 \(S_i^{n_i}\) 的映射在各副本间的坐标映射都为零。因此每个自同态保持每个 \(S_i^{n_i}\)，不同类型之间没有非对角块，而各块内的映射可以独立选取，得到
\[
\operatorname{End}_R({}_RR)
\cong\prod_i\operatorname{End}_R(S_i^{n_i})
\cong\prod_iM_{n_i}(E_i).
\]
第二个同构是在固定各副本与 \(S_i\) 的识别后，把自同态写成矩阵，条目表示“从第 \(j\) 个副本到第 \(i\) 个副本”的映射。若两个自同态的条目分别为 \(f_{ij},g_{ij}\)，先施加 \(g\)、再施加 \(f\)，从第 \(\ell\) 个副本到第 \(i\) 个副本的条目为 \(\sum_jf_{ij}\circ g_{j\ell}\)。这恰好是以复合作为 \(E_i\) 乘法的通常矩阵乘法，也解释了为何同一简单类型的副本形成一个矩阵块。

由\cref{b2:prop:regular-hom-endomorphism}，每个左线性自同态是某个右乘映射 \(r\mapsto ra\)；先右乘 \(b\)、再右乘 \(a\)，结果为右乘 \(ba\)。因此
\[
R^{\mathrm{op}}\cong\operatorname{End}_R({}_RR)
\cong\prod_iM_{n_i}(E_i).
\]
要回到原环 \(R\)，先对整个式子取反环，得到
\[
R\cong\prod_iM_{n_i}(E_i)^{\mathrm{op}}.
\]
这里得到的仍是矩阵环的反环。对任意除环 \(E\) 和 \(n\ge1\)，用转置把它写回通常的矩阵形式：
\[
\Phi:M_n(E)^{\mathrm{op}}\longrightarrow M_n(E^{\mathrm{op}}),
\qquad A^{\mathrm{op}}\longmapsto A^{\mathsf T}.
\]
来源反转整个矩阵的乘法，目标则反转每个条目的乘法；转置同时交换行列，才能把求和中的索引对应起来。确实，设 \(A=(a_{ij})\)、\(B=(b_{ij})\) 属于 \(M_n(E)\)，来源中 \(A^{\mathrm{op}}B^{\mathrm{op}}=(BA)^{\mathrm{op}}\)，目标中的乘积逐条目满足
\[
\bigl(A^{\mathsf T}\cdot_{E^{\mathrm{op}}}B^{\mathsf T}\bigr)_{i\ell}
=\sum_j a_{ji}\cdot_{\mathrm{op}}b_{\ell j}
=\sum_j b_{\ell j}a_{ji}
=(BA)_{\ell i}.
\]
这正是 \((BA)^{\mathsf T}\) 的第 \((i,\ell)\) 项。转置还保持加法与单位，并且再转置就恢复原矩阵，所以 \(\Phi\) 为含幺环同构。取 \(D_i=E_i^{\mathrm{op}}\)，便得到题述的原环分解。

反过来，设 \(D\) 为除环、\(n\ge1\)、\(A=M_n(D)\)。右乘矩阵单位 \(E_{jj}\) 保留第 \(j\) 列而清去其余列，因此按列得到左模直和
\[
{}_AA=AE_{11}\oplus\cdots\oplus AE_{nn}.
\]
每个左理想 \(AE_{jj}\) 由只有第 \(j\) 列可能非零的矩阵组成，提取该列便给出它与列模 \(D^n\) 的左 \(A\) 模同构。以标准列为基，\(D^n\) 是系数写在右侧的右 \(D\) 向量空间，而 \(A\) 通过矩阵乘法从左作用。这个左 \(A\) 模简单：若 \(v\ne0\)，选 \(v_j\ne0\)；要把 \(v\) 送到任意列向量 \(w\)，可读取第 \(j\) 坐标，用其逆消去 \(v_j\)，再把各 \(w_i\) 写入目标坐标。具体取只有第 \(j\) 列可能非零的矩阵 \(B\)，令 \(b_{ij}=w_i v_j^{-1}\)，则
\[
(Bv)_i=(w_i v_j^{-1})v_j=w_i.
\]
因此任何非零向量都生成全模。乘法次序不能改成 \(v_j^{-1}w_i\)，因为除环不必交换。故每个矩阵环左半单；有限乘积的正则模按各因子分解，仍为有限个简单左模的直和，其余因子在该分量上作用为零。

矩阵环的右正则模则按行分解为 \(E_{11}A\oplus\cdots\oplus E_{nn}A\)。对非零行向量 \(v\)，取 \(v_j\ne0\)；给定目标行 \(w\)，令 \(B\) 只有第 \(j\) 行可能非零，且 \(b_{j\ell}=v_j^{-1}w_\ell\)，便有 \((vB)_\ell=v_j(v_j^{-1}w_\ell)=w_\ell\)。所以各行理想是简单右模，有限乘积也右半单。两侧正则模都只有有限个简单分量，因而具有有限长度；分别把它们看作左 \(R\) 模与左 \(R^{\mathrm{op}}\) 模，\cref{b2:thm:finite-length-chain-conditions}给出两侧的升链与降链条件。零环的情形各项结论直接成立。
\end{proof}

这一定理称为\term[Artin-Wedderburn-dingli]{Artin–Wedderburn 定理}[Artin--Wedderburn theorem]。它把所有半单环归结为除环及其有限矩阵环。这些除环可以非交换；有限维的实代数就可能含有四元数除环这一类型。此前半单环只按左正则模定义，右半单性是在得到矩阵结构后由行分解证明的结论。

\subsection{单 Artin 环与除环}
\label{b2:subsec:1303:02}

\begin{definition}\label{b2:def:simple-ring}
设 \(R\) 为非零含幺环。若 \(R\) 只有 \(0,R\) 两个双边理想，就称为\term[jiandan-huan]{单环}[simple ring]。它与左正则模 \({}_RR\) 简单是不同条件：后者要求只有 \(0,R\) 两个左理想，因而更强。
\end{definition}

\begin{proposition}\label{b2:prop:simple-artinian}
非零含幺环 \(R\) 是单环且左 Artin，当且仅当 \(R\) 同构于某个 \(M_n(D)\)，其中 \(D\) 为除环、\(n\ge1\)。这些环也右 Artin。特别地，单半单环恰为单个除环上的全矩阵环。
\end{proposition}
\begin{proof}
先设 \(A=M_n(D)\)。上一定理已证明 Artin 性。若双边理想 \(I\) 含非零矩阵 \(B\)，取 \(b_{ij}\ne0\)。左右乘矩阵单位可以把这个非零条目搬到任意位置 \((k,\ell)\)：
\[
E_{ki}BE_{j\ell}=b_{ij}E_{k\ell}\in I.
\]
再左乘对角矩阵 \(b_{ij}^{-1}I_n\)，把该条目化成一，得到 \(E_{k\ell}\in I\)。对每个 \(a\in D\)，继续左乘 \(aI_n\)，还得到 \(aE_{k\ell}\in I\)。任意矩阵是这些矩阵的有限和，所以 \(I=A\)，证明单性。

反过来，设 \(R\) 单且左 Artin。非零左理想的降链条件保证存在极小非零左理想，它作为左模简单。要利用单性，我们还需要由这些简单左理想构造一个非零双边理想，所以令 \(U\) 为全部极小左理想之和。它首先是非零左理想。对简单左理想 \(L\) 和 \(r\in R\)，右乘映射
\[
\rho_r:L\longrightarrow R,\qquad x\longmapsto xr
\]
为左 \(R\) 线性，因为 \((ax)r=a(xr)\)。它的核是 \(L\) 的子模，简单性使核为零或全体 \(L\)；因此像 \(Lr\) 或者为零，或者由第一同构定理同构于 \(L\)，仍是简单左理想。故 \(Lr\subseteq U\)，对 \(U\) 中的有限和也有 \(Ur\subseteq U\)。于是 \(U\) 为非零双边理想，由单性 \(U=R\)。\cref{b2:thm:semisimple-module-characterizations}说明 \({}_RR\) 半单，再由\cref{b2:thm:artin-wedderburn}得到有限矩阵环乘积。若乘积有两个以上非零因子，其中一个因子便给出非零真双边理想，与单性矛盾。因此只剩一个因子。
\end{proof}

例如 \(M_n(D)\) 在 \(n>1\) 时是单环，却不是除环：\(E_{11}\) 非零而不可逆，且 \(AE_{11}\) 是非零真左理想。相反，若非零含幺环 \(R\) 的左正则模本身简单，对任意 \(0\ne a\in R\) 有 \(Ra=R\)，于是存在 \(b\) 使 \(ba=1\)。再对 \(b\ne0\) 取 \(c\) 使 \(cb=1\)，则 \(c=cba=a\)，所以 \(ab=1\)。因此每个非零元都有双侧逆，\(R\) 是除环。反过来，除环的每个非零左理想都含有单位，因而为全环，所以其左正则模简单。

\ProseParagraphBreak
上述证明还说明：在单环中，只要存在一个极小非零左理想，就已经可以推出矩阵环结构；左 Artin 条件提供了保证其存在的充分条件。随后用单性把全部极小左理想的和扩大为全环，环的双边理想结构与左模结构分别承担了不同步骤。

\subsection{分解的唯一性及左右约定}
\label{b2:subsec:1303:03}

一个半单分解可以改变简单分量的具体位置，却不应改变每种简单类型的副本数。固定简单模 \(S\)，从 \(S\) 到其他简单类型的同态为零，而到一个 \(S\) 副本的同态给出一个自同态除环。因此 \(\operatorname{Hom}_R(S,M)\) 可以检测这种类型的副本数。其标量来自来源 \(S\) 的自同态，作用方式是先改变输入，再施加同态，即预复合。

\begin{proposition}[半单分解中的重数]\label{b2:prop:semisimple-multiplicity}
设 \(R\) 为含幺环、\(M\) 为有限生成半单幺左 \(R\) 模、\(S\) 为简单幺左 \(R\) 模，\(E=\operatorname{End}_R(S)\)，其乘法为映射复合。则 \(\operatorname{Hom}_R(S,M)\) 通过以下运算成为右 \(E\) 向量空间（\(h\in\operatorname{Hom}_R(S,M),\ e\in E\)）：
\[
h\cdot e=h\circ e.
\]
任一简单直和分解中，\(S\) 的同构副本出现次数等于以下右维数：
\[
\dim_E\operatorname{Hom}_R(S,M).
\]
特别地，简单类型及各自重数不依赖所选分解。
\end{proposition}
\begin{proof}
由复合的结合律，
\[
(h\cdot e)\cdot e'=h\circ e\circ e'=h\cdot(ee'),
\]
而恒等映射给出单位作用，复合对加法分配，所以确为右作用。取简单分解 \(M=\bigoplus_{j\in J}T_j\)。由于 \(S\) 简单，任意 \(0\ne s\in S\) 都生成 \(S\)。给定 \(h:S\to M\)，\(h(s)\) 只含有限多个非零坐标，左标量作用不会增加这些坐标；由 \(h(rs)=rh(s)\)，整个 \(h(S)\) 都在同一个有限子直和中。因此各坐标映射组成有限支集的族，取坐标投影给出右 \(E\) 线性同构
\[
\operatorname{Hom}_R(S,M)\cong\bigoplus_{j\in J}\operatorname{Hom}_R(S,T_j).
\]
反向把有限多个坐标映射通过包含求和即可。这里即使 \(J\) 无限，得到的也是直和，而不是直积，原因在于来源 \(S\) 有限生成。又由于 \(M\) 有限生成，由\secref{b2:sec:1301}的有限支集论证知 \(J\) 有限。若 \(T_j\not\cong S\)，该项由 Schur 引理为零；若 \(u_j:S\to T_j\) 是同构，则每个同态唯一写为 \(u_j\circ e\)，所以该项是一维右 \(E\) 空间。维数恰为副本数。反过来，\(S\) 出现当且仅当这个 Hom 非零，因此出现的简单类型也由 \(M\) 本身决定。
\end{proof}

这里除环上的有限维数仍然唯一。基交换论证只需要非零系数可逆：若右线性展开含 \(v_j a\)、\(a\ne0\)，就在右侧乘 \(a^{-1}\) 解出 \(v_j\)，然后逐个替换生成元。一组无关向量不能多于一组有限生成元，交换两组基便得大小相同；不需要交换两个标量的次序。

\ProseParagraphBreak
将上一命题用于半单环的正则模，\(\operatorname{Hom}_R(S,R)\) 的右维数便恢复类型 \(S\) 的重数，而该类型还决定系数除环
\[
D_S=\operatorname{End}_R(S)^{\mathrm{op}}.
\]
这两种数据取自 \(R\) 的模与同态，不依赖矩阵坐标。不同矩阵因子则可以由环内部的中心幂等元区分。

\begin{proposition}\label{b2:prop:matrix-centers-central-idempotents}
设 \(D\) 为除环、\(n\ge1\)，则
\[
Z\bigl(M_n(D)\bigr)=Z(D)I_n.
\]
更一般地，设 \(t\ge0\)、\(D_1,\ldots,D_t\) 为除环、\(n_1,\ldots,n_t\ge1\)，并令 \(R=\prod_{i=1}^tM_{n_i}(D_i)\)，则
\[
Z(R)=\prod_{i=1}^t Z(D_i)I_{n_i}.
\]
对 \(1\le i\le t\)，令 \(c_i\in R\) 的第 \(i\) 坐标为 \(I_{n_i}\)，其余坐标为零。\(R\) 的中心幂等元恰为 \(\sum_{i\in F}c_i\)，其中 \(F\subseteq\{1,\ldots,t\}\)，总数为 \(2^t\)。按 \(e\le f\iff ef=e\) 比较中心幂等元时，极小非零者恰为 \(c_1,\ldots,c_t\)。\(t=0\) 时 \(R\) 为零环，只有幂等元零，没有非零中心幂等元。
\end{proposition}
\begin{proof}
设 \(C=(c_{ij})\in Z(M_n(D))\)。与各 \(E_{jj}\) 交换，迫使第 \(j\) 行、第 \(j\) 列的非对角元都为零，因此 \(C\) 为对角矩阵。再与每个 \(E_{ij}\) 交换，得到各对角元相同，故 \(C=dI_n\)。对任意 \(a\in D\)，与 \(aE_{11}\) 交换又给出 \(da=ad\)，所以 \(d\in Z(D)\)。反过来，\(d\in Z(D)\) 时，\(dI_n\) 与每个矩阵逐条目交换，证明第一个等式。这些论证在 \(n=1\) 时同样成立。

乘积环中的元素与全部元素交换，当且仅当每个坐标在相应因子的中心，得到第二个等式。\(Z(D_i)\) 是域：非零中心元的逆仍与每个元素交换。故其中的幂等元只能为零或一，中心幂等元便是在各坐标上分别选择零或单位矩阵，恰对应子集 \(F\)，共有 \(2^t\) 个。两个这样的幂等元相乘对应子集的交，所以 \(e\le f\) 对应子集包含；极小非空子集是单点集，给出各 \(c_i\)。空乘积时仅有空子集，对应零环中的唯一元素零，仍有 \(2^0=1\) 个幂等元。
\end{proof}

\begin{corollary}\label{b2:cor:artin-wedderburn-uniqueness}
设 \(R\) 为含幺环，且有两个含幺环分解
\[
R\cong\prod_{i=1}^tM_{n_i}(D_i)
\cong\prod_{j=1}^sM_{m_j}(F_j),
\]
其中 \(t,s\ge0\)，\(D_i,F_j\) 为除环，\(n_i,m_j\ge1\)。则 \(t=s\)，并可重排下标，使每个 \(i\) 都满足 \(n_i=m_i\) 及 \(D_i\cong F_i\)（含幺环同构）。零环的两个分解均为空乘积。
\end{corollary}
\begin{proof}
若 \(R=0\)，任何除环上的正阶矩阵环都非零，所以只能有 \(t=s=0\)。以下设 \(R\ne0\)。把两种乘积的各坐标单位通过同构送到 \(R\)，分别得到 \(c_1,\ldots,c_t\) 与 \(d_1,\ldots,d_s\)。由\cref{b2:prop:matrix-centers-central-idempotents}，它们分别列出了 \(R\) 的全部极小非零中心幂等元。因此 \(t=s\)，且重排后 \(c_i=d_i\)；各因子正是这些幂等元切出的环 \(c_iR\)，其单位为 \(c_i\)。

对任意幺左 \(R\) 模 \(T\)，各 \(c_iT\) 是子模，因为 \(c_i\) 位于中心。由 \(\sum_i c_i=1\) 及 \(c_ic_j=0\ (i\ne j)\)，有
\[
T=\bigoplus_{i=1}^t c_iT.
\]
若 \(T\) 简单，就只能有一个非零分量，该分量等于 \(T\)，其余因子在 \(T\) 上作用为零。因此不同矩阵因子的简单类型不会混合。对 \(A=M_n(D)\)，前面的按列分解给出 \({}_AA\cong(D^n)^n\)。每个简单左 \(A\) 模 \(T\) 都由任意非零元生成，因而是 \(A\) 的商；相应满射 \(A\to T\) 在某个列理想上的限制非零，Schur 引理使 \(T\cong D^n\)。故一个矩阵因子恰提供一种简单类型。

还要从这个类型恢复 \(D\)。令 \(e_i\) 为 \(D^n\) 的标准列，\(f\in\operatorname{End}_A(D^n)\)。由 \(E_{11}f(e_1)=f(E_{11}e_1)=f(e_1)\)，有 \(f(e_1)=e_1a\)，其中 \(a\in D\)。对任意列 \(v\)，取第一列为 \(v\)、其余列为零的矩阵 \(B\)，则 \(Be_1=v\)，左 \(A\) 线性给出
\[
f(v)=f(Be_1)=Bf(e_1)=B(e_1a)=va.
\]
反过来，每个右乘映射 \(r_a(v)=va\) 都左 \(A\) 线性，因为 \(B(va)=(Bv)a\)。又 \(r_a\circ r_b=r_{ba}\)，故
\[
\operatorname{End}_A(D^n)\cong D^{\mathrm{op}}.
\]
于是该简单类型的自同态除环再取反环，恰好恢复原来的 \(D\)。

对因子 \(c_iR\)，两种分解给出同一个简单左 \(R\) 模类型 \(S_i\)；因其余因子作用为零，左 \(R\) 线性与左 \(c_iR\) 线性在 \(S_i\) 上是同一条件。因此其自同态环在第一种描述下同构于 \(D_i^{\mathrm{op}}\)，在第二种描述下同构于 \(F_i^{\mathrm{op}}\)，取反环得到 \(D_i\cong F_i\)。按列分解又说明 \(S_i\) 在正则模中的重数分别为 \(n_i\) 与 \(m_i\)。\({}_RR\) 有限生成且半单，\cref{b2:prop:semisimple-multiplicity}使这两个重数都等于 \(\operatorname{Hom}_R(S_i,R)\) 的右 \(\operatorname{End}_R(S_i)\) 维数，因此 \(n_i=m_i\)。
\end{proof}

这里所谓唯一性是参数及因子的同构类型唯一，不是矩阵单位或具体同构映射唯一。改变同一简单类型的副本识别，会改变矩阵坐标。采用左模约定时，自同态除环必须再取反环才得到原环的矩阵系数；若采用右模约定，相应次序随之改变，不能在公式中无说明地省去 \(\mathrm{op}\)。

% !TeX root = ../../Book2.tex
% 纲要标题：### 13.4 有限维代数、群代数与有限环
% 纲要位置：计划纲要.md，第 1482 行。

\section{有限维代数、群代数与有限环}
\label{b2:sec:1304}

结构定理不仅给出抽象的环同构，也使维数、简单模和标量域之间的关系可以计算。本节先讨论有限维代数，再用最小的群代数例子比较不同特征，最后证明有限除环必交换。最后一个结论需要环本身只有有限多个元素，不能仅凭它在某个域上有限维就套用。

% 纲要要点（供目录核对）：
%
% * 有限维半单代数
% * 群代数的半单性例子与 Maschke 定理的陈述
% * 有限除环的交换性与 Wedderburn 小定理；循环模的无重因子判据
%
% ---
%

\subsection{有限维半单代数}
\label{b2:subsec:1304:01}

\begin{proposition}\label{b2:prop:finite-dimensional-semisimple-algebra}
设 \(k\) 为域、\(A\) 为有限维结合含幺半单 \(k\) 代数。存在非负整数 \(t\)、有限维除 \(k\) 代数 \(D_1,\ldots,D_t\) 及正整数 \(n_1,\ldots,n_t\)，使
\[
A\cong\prod_{i=1}^tM_{n_i}(D_i)
\]
成为 \(k\) 代数同构。相应的简单左模为列模 \(S_i=D_i^{n_i}\)，由第 \(i\) 个因子作用，其余因子作用为零，且
\[
\dim_kA=\sum_i n_i^2\dim_kD_i,
\qquad \dim_kS_i=n_i\dim_kD_i.
\]
若 \(k\) 代数闭，则所有 \(D_i=k\)，于是 \(A\cong\prod_iM_{n_i}(k)\)，且 \(\dim_kA=\sum_i(\dim_kS_i)^2\)。零代数对应 \(t=0\) 的空乘积，此时简单模列表为空，维数公式中的空和为零。
\end{proposition}
\begin{proof}
零代数的结论由空乘积约定成立。设 \(A\ne0\)。对任意简单左 \(A\) 模 \(S_i\)，取 \(0\ne s_i\in S_i\)，则 \(As_i=S_i\)，因为这是非零子模。因此 \(a\mapsto as_i\) 给出 \(k\) 线性满射 \(A\to S_i\)，从而 \(\dim_kS_i\le\dim_kA<\infty\)。\(\operatorname{End}_A(S_i)\) 是 \(\operatorname{End}_k(S_i)\) 的 \(k\) 子空间，故有限维；Schur 引理使它为除环，其反环 \(D_i\) 也是有限维除 \(k\) 代数。\cref{b2:thm:artin-wedderburn}的构造保持 \(k\) 标量：标量在正则模上作为乘 \(\lambda\) 作用，在各简单分量上仍是同一个标量自同态，取反环不会改变中心标量。因此得到 \(k\) 代数同构。矩阵有 \(n_i^2\) 个任意 \(D_i\) 条目，列模有 \(n_i\) 个条目，直接得到两条维数公式。若 \(k\) 代数闭，刚证明的简单模有限维性允许应用\cref{b2:cor:schur-algebraically-closed}，使各自同态除环等于 \(k\)，反环仍为 \(k\)，从而得到平方和公式。
\end{proof}

例如，代数闭域上的一个四维半单代数只能同构于 \(k^4\) 或 \(M_2(k)\)：其维数是正整数平方之和，只有 \(4=1+1+1+1\) 与 \(4=2^2\) 两种可能。前者有四种一维简单模，后者只有一种二维简单模。两个代数的总维数相同，模结构却不同。

\ProseParagraphBreak
一般基域上不能删去除代数的维数。比如 \(\mathbb C\) 作为实代数是半单的，其唯一简单模为 \(\mathbb C\)，实维数为二；公式是 \(\dim_{\mathbb R}\mathbb C=1^2\cdot2\)，而不是这个简单模实维数的平方。类似地，实四元数除代数维数为四，它的左正则模简单；有限维除代数仍可能非交换。

\subsection{群代数的半单性}
\label{b2:subsec:1304:02}

设 \(k\) 为域、\(G\) 为有限群。群代数 \(k[G]\) 以各 \(g\in G\) 为 \(k\) 基，乘法按群乘法线性延拓。第五卷第1.3节将通过对群作用平均投影，证明 Maschke 定理\footnote{马施克（Heinrich Maschke，1853--1908），德国数学家，研究有限群的线性表示和微分几何。}：\(k[G]\) 半单，当且仅当 \(|G|\cdot1_k\ne0\)，也就是 \(|G|\cdot1_k\) 在 \(k\) 中可逆。

\ProseParagraphBreak
平均要修正的是普通线性投影与群作用不相容的问题。若 \(\rho:G\to\operatorname{GL}_k(V)\) 是表示，\(U\subseteq V\) 在群作用下稳定，先选一个线性投影 \(p:V\to U\)，并通过包含把它看作 \(V\) 的自同态。在 \(|G|\cdot1_k\ne0\) 时，平均得到
\[
P=\frac1{|G|}\sum_{g\in G}\rho(g)p\rho(g)^{-1}.
\]
群作用会排列这些共轭项，使平均后的算子与群作用相容。每个共轭项在 \(U\) 上都是恒等，所以未经归一化的和在 \(U\) 上是 \(|G|\) 倍的恒等；只有能除以 \(|G|\)，才恢复所需的投影。这解释了特征零或 \(\operatorname{char}k=p>0\) 且 \(p\mathrel{\vcenter{\hbox{\scalebox{1.25}{\(\nmid\)}}}}|G|\) 的条件。

\ProseParagraphBreak
最小例子可以在当前知识范围内直接算出。设 \(G=C_2=\langle g\mid g^2=1\rangle\)，且 \(\operatorname{char}k\ne2\)。令
\[
e_+=\frac{1+g}{2},\qquad e_-=\frac{1-g}{2}.
\]
由 \(g^2=1\) 得 \(e_+^2=e_+\)、\(e_-^2=e_-\)、\(e_+e_-=0\)、\(e_++e_-=1\)，并且 \(ge_+=e_+\)、\(ge_-=-e_-\)。因此 \(k[C_2]=ke_+\oplus ke_-\)，左乘 \(e_+\) 与 \(e_-\) 分别是正则作用中 \(g\) 的 \(1\) 与 \(-1\) 特征子空间上的投影。这就是\cref{b2:prop:coprime-operator-decomposition}中用互素多项式构造投影的情形：\(t-1\) 与 \(t+1\) 互素，\((1+t)/2\) 与 \((1-t)/2\) 分别选出两个分量。具体地，
\[
k[C_2]\longrightarrow k\times k,
\quad a+bg\longmapsto(a+b,a-b)
\]
为代数同构，逆映射为 \((u,v)\mapsto u e_++v e_-\)。两种简单模分别使 \(g\) 作用为 \(1\) 与 \(-1\)，正则模就是这两个一维模的直和。

\ProseParagraphBreak
若 \(\operatorname{char}k=2\)，令 \(\varepsilon=g-1\)，则 \(\varepsilon^2=g^2-2g+1=0\)，从而
\[
k[C_2]\cong k[\varepsilon]/(\varepsilon^2).
\]
这里 \(\varepsilon\ne0\)，因为 \(1,g\) 仍是群代数的两个基向量。两个特征值 \(1,-1\) 在特征二时合并，留下非零的幂零元 \(\varepsilon\)。子模 \(k\varepsilon\) 与相应的商模都为一维，\(g\) 均作用为恒等，但这个子模没有直和补模；\cref{b2:prop:semisimple-cyclic-polynomial}将通过重复因子的一般判据证明这种不分裂，也给出该环不半单的结论。相同的群在不同特征下给出了不同的模分裂行为。

\ProseParagraphBreak
对 \(G=C_m\)、\(m\ge1\)，\cref{b2:prop:semisimple-cyclic-polynomial}的证明将给出同构 \(k[C_m]\cong k[t]/(t^m-1)\)，其中 \(t\) 的陪集对应循环群生成元 \(g\)。因而一个 \(C_m\) 表示就是一个满足 \(T^m=\operatorname{id}\) 的算子，群代数的模也正是由这个关系限制的 \(k[t]\) 模。这把问题接回第四章的算子模与主分量：不同不可约因子将分量分开，重复因子则可能留下不能分裂的幂零层。循环群的半单性也将在该命题中判定；一般有限群仍由第五卷的平均方法处理。

\subsection{Wedderburn 小定理}
\label{b2:subsec:1304:03}

\begin{theorem}[Wedderburn 小定理]\label{b2:thm:wedderburn-little}
每个只有有限多个元素的除环都是交换域。
\end{theorem}
\begin{proof}
设 \(D\) 为有限除环。中心 \(Z\) 是域：非零中心元 \(z\) 的逆仍与每个元素交换，因为由 \(zx=xz\) 左右乘以 \(z^{-1}\) 就得到 \(xz^{-1}=z^{-1}x\)。记 \(|Z|=q\ge2\)，把 \(D\) 视为有限维 \(Z\) 向量空间，令 \(n=\dim_ZD\)，则 \(|D|=q^n\)。若 \(n=1\)，每个元素都是中心中 \(1\) 的倍数，结论成立。下面假设 \(n>1\)，推出矛盾。

乘法群 \(D^\times\) 的中心为 \(Z^\times\)，含 \(q-1\) 个元素。为了计算非中心共轭类的大小，先考察代表元的中心化子。对 \(x\in D\setminus Z\)，其环中心化子
\[
C_D(x)=\{\,a\in D\mid ax=xa\,\}
\]
是包含 \(Z\) 的除子环：加、乘封闭，若 \(a\ne0\) 与 \(x\) 交换，则其逆也交换。设 \(m_x=\dim_Z C_D(x)\)，则 \(|C_D(x)|=q^{m_x}\)。虽然 \(C_D(x)\) 未必交换，\(D\) 仍是它的左向量空间，而 \(Z\) 位于 \(C_D(x)\) 的中心。若 \(c_1,\ldots,c_{m_x}\) 是 \(C_D(x)\) 的 \(Z\) 基，\(v_1,\ldots,v_{d_x}\) 是 \(D\) 的左 \(C_D(x)\) 基，则 \(c_i v_j\) 是 \(D\) 的 \(Z\) 基：先按 \(v_j\) 展开，再把左侧的系数按 \(c_i\) 展开，且两次展开都唯一。因此维数塔公式给出
\[
\dim_ZD=\dim_{C_D(x)}D\cdot\dim_ZC_D(x),
\qquad n=d_xm_x.
\]
这里的乘法次序始终是左系数乘基向量，证明不需要 \(C_D(x)\) 交换。故 \(m_x\mid n\)；又因 \(x\) 非中心，\(C_D(x)\) 是 \(D\) 的真 \(Z\) 子空间，所以 \(m_x<n\)。

群中心化子只取其中的可逆元：
\[
C_{D^\times}(x)=\{a\in D^\times\mid ax=xa\}=C_D(x)^\times,
\]
所以它的阶为 \(q^{m_x}-1\)。有限群的类方程把中心的单元素共轭类与其余共轭类分开，写成 \(|G|=|Z(G)|+\sum_j[G:C_G(x_j)]\)，其中 \(x_j\) 取遍非中心共轭类的代表。把第一卷命题4.2.2及式(4.1)的这一公式应用于 \(D^\times\)，取代表 \(x_1,\ldots,x_s\)，便得
\begin{equation}\label{b2:eq:wedderburn-class-equation}
q^n-1=(q-1)+\sum_{j=1}^s\frac{q^n-1}{q^{m_{x_j}}-1}.
\end{equation}
每个分数都是该代表的共轭类大小，而各分母的指数 \(m_{x_j}\) 都是 \(n\) 的真因子。

要从这个整数等式中得到约束，可以寻找一个同时整除左端与所有非中心共轭类大小的整数。分母的指数都是 \(n\) 的真因子，使第一卷定义21.1.1中的分圆因子 \(\Phi_n\) 正好具有这一性质。该卷式(21.1)及其后的整数系数归纳证明给出
\[
t^n-1=\prod_{d\mid n}\Phi_d(t),\qquad \Phi_d(t)\in\mathbb Z[t].
\]
对每个真因子 \(m\mid n\)，有整数系数多项式恒等式
\[
\frac{t^n-1}{t^m-1}=\prod_{\substack{d\mid n\\d\nmid m}}\Phi_d(t).
\]
因为 \(m<n\)，\(n\) 不整除 \(m\)，右端包含 \(\Phi_n(t)\)。代入 \(t=q\)，便知整数 \(\Phi_n(q)\) 整除 \(q^n-1\)，也整除式\eqref{b2:eq:wedderburn-class-equation}中的每个分数。从左端减去这些均被它整除的共轭类大小，剩下的正是 \(q-1\)，所以
\[
\Phi_n(q)\mid q-1.
\]

但 \(\Phi_n(q)>q-1\)。此时 \(\Phi_n(q)\) 已是一个整数；把同一个多项式放到复数域中分解，只是为了估计这个整数的大小。按其定义写
\[
\Phi_n(q)=\prod_{\zeta\text{ 为本原 }n\text{ 次单位根}}(q-\zeta).
\]
因 \(n>1\)，每个这样的 \(\zeta\ne1\)。单位圆上除 \(1\) 外的点都比 \(1\) 离正实数 \(q>1\) 更远，具体地，
\[
|q-\zeta|^2=(q-1)^2+2q(1-\operatorname{Re}\zeta)>(q-1)^2.
\]
非实根成共轭对，每对贡献 \((q-\zeta)(q-\overline\zeta)=|q-\zeta|^2>0\)；若有实本原根，则只能是 \(n=2\) 时的 \(-1\)，因子 \(q+1\) 也为正。因此 \(\Phi_n(q)>0\)，且它等于上述所有因子绝对值的乘积，严格大于 \((q-1)^{\varphi(n)}\ge q-1\)。这里 \(\varphi(n)\ge1\)，而 \(q-1\ge1\)；即使 \(q=2\)，每个因子绝对值也严格大于 \(1\)，严格不等式仍成立。这与正整数 \(\Phi_n(q)\) 整除 \(q-1\) 矛盾。故 \(n=1\)，\(D=Z\) 交换。
\end{proof}

上述结论称为\term[Wedderburn-xiaodingli]{Wedderburn 小定理}[Wedderburn's little theorem]。分圆多项式的整数性及因式分解给出关键的整除关系；有限性使中心和各中心化子的元素个数分别写成 \(q\) 的整数次幂，也使乘法群的类方程成为有限整数等式。本定理要求除环只有有限多个元素；例如四元数除代数 \(\mathbb H\) 虽在 \(\mathbb R\) 上四维，却有无限多个元素，并不交换。

\begin{corollary}\label{b2:cor:finite-semisimple-rings}
设 \(R\) 为含幺环。\(R\) 有限且半单，当且仅当存在非负整数 \(t\)、正整数 \(n_1,\ldots,n_t\) 及素数幂 \(q_1,\ldots,q_t\)，使
\[
R\cong\prod_{i=1}^t M_{n_i}(\mathbb F_{q_i}).
\]
\(t=0\) 时把空乘积解释为零环。
\end{corollary}
\begin{proof}
由\cref{b2:thm:artin-wedderburn}分解后，每个矩阵因子仍有限，其系数除环作为单个矩阵条目的取值集也有限。\cref{b2:thm:wedderburn-little}使每个除环都是有限域。反过来，有限个有限域上的有限阶矩阵环显然只有有限多个元素，且由结构定理半单。
\end{proof}

回到循环群代数，交换性使半单分解中的矩阵因子只能是一阶的域。对多项式商环，这种分解可以直接从不可约因子读出；若因子重复，问题则是相应的子模能否真正分裂，而不只是有几个简单的合成因子。

\begin{proposition}\label{b2:prop:semisimple-cyclic-polynomial}
设 \(k\) 为域，\(f\in k[t]\) 为非零非常数多项式。以下条件等价：\(k[t]/(f)\) 作为幺左 \(k[t]\) 模半单；它作为含幺环半单；\(f\) 的每个不可约因子在 \(k[t]\) 中只出现一次。

若 \(m\ge1\) 为整数、\(C_m\) 为 \(m\) 阶循环群，则群代数 \(k[C_m]\) 半单，当且仅当 \(\operatorname{char}k=0\)，或 \(\operatorname{char}k=p>0\) 且 \(p\) 不整除 \(m\)。
\end{proposition}
\begin{proof}
令 \(A=k[t]/(f)\)。因为 \(k[t]\to A\) 满射，\(A\) 的左 \(k[t]\) 子模与左 \(A\) 子模相同，简单性和直和分解也相同。因此 \(A\) 作为 \(k[t]\) 模半单，等价于左正则模 \({}_AA\) 半单，也就是环 \(A\) 半单。

若 \(f=c\,p_1\cdots p_r\)，其中 \(c\in k^\times\)，\(p_i\) 是两两不同的首一不可约多项式，则这些因子两两互素。第一卷定理10.4.3的中国剩余定理给出
\[
A\cong\prod_{i=1}^r k[t]/(p_i).
\]
每个因子都是域，也是简单 \(k[t]\) 模；有限乘积作为模就是有限直和，故 \(A\) 半单。

若某个不可约多项式 \(p\) 满足 \(p^2\mid f\)，则 \(A\) 有商模 \(Q=k[t]/(p^2)\)。设 \(A\) 半单，\cref{b2:prop:semisimple-subquotients}便使 \(Q\) 半单，故非零子模 \(pQ\) 应有直和补模，相应地存在 \(k[t]\) 线性投影 \(u:Q\to pQ\)，在 \(pQ\) 上为恒等。记 \(1,p\) 在 \(Q\) 中的陪集为 \(\overline1,\overline p\)。因 \(p^2Q=0\)，\(p\) 消去 \(pQ\)，而 \(u(\overline1)\in pQ\)，于是
\[
u(\overline p)=u(p\overline1)=p\,u(\overline1)=0.
\]
但 \(\overline p\ne0\)，且 \(\overline p\in pQ\)，所以投影又要求 \(u(\overline p)=\overline p\)，矛盾。故 \(A\) 不半单。

取 \(C_m=\langle g\rangle\)，映射 \(t\mapsto g\) 给出代数同构 \(k[t]/(t^m-1)\cong k[C_m]\)，因为它把前者的基 \(\overline1,\overline t,\ldots,\overline{t^{m-1}}\) 映到后者的基 \(1,g,\ldots,g^{m-1}\)。若 \(\operatorname{char}k=0\)，或特征 \(p\) 不整除 \(m\)，则 \(m\cdot1_k\ne0\)，而
\[
(t^m-1)'=mt^{m-1},
\qquad \gcd(t^m-1,mt^{m-1})=1.
\]
任何重复的不可约因子都会同时整除多项式及其导数，因此 \(t^m-1\) 没有重复不可约因子，群代数半单。若 \(\operatorname{char}k=p>0\) 且 \(p\mid m\)，则
\[
t^m-1=(t^{m/p}-1)^p.
\]
括号内仍为非常数多项式，至少有一个不可约因子，而它在右端的重数至少为 \(p\ge2\)，所以群代数不半单。
\end{proof}


\begin{chaptersummary}
半单模的关键性质是每个子模都有补模。简单子模之和可以抽成直和，半单模的子模与商模仍半单；但两端半单的扩张未必半单。Schur 引理把不同简单类型间的同态消去，并把同一简单类型的自同态组成除环。

\ProseParagraphBreak
半单环的正则模由有限多个简单左理想组成，这使整个左模范畴中的短正合列都分裂。把正则模的自同态按简单类型分块，再通过反环恢复原乘法，得到有限个矩阵环的乘积。简单类型决定系数除环，正则模中的重数决定矩阵大小；具体坐标可以改变，这些参数的同构类型不变。

\ProseParagraphBreak
半本原性只要求所有简单模共同零化的元素为零，它本身不保证正则模半单。有限维代数的结构还受基域影响，有限除环的交换性则依靠更强的“元素总数有限”。Wedderburn 小定理把有限半单环进一步归结为有限域上的矩阵环；下一章将用 Jacobson 根研究一般环偏离半单结构的部分。
\end{chaptersummary}

\BookExercises

\begin{exercise}\label{b2:ex:13-simple-polynomial-modules}
设 \(k\) 为域。分类 \(k[t]\) 的全部幺性简单左模，并说明哪些具有 \(k\) 维数一。在 \(k=\mathbb R\) 时给出一个二维简单模。
\end{exercise}
\begin{solution}
每个简单模是极大理想的商。\(k[t]\) 是主理想整环，极大理想恰为不可约多项式 \(f\) 生成的理想，所以简单模为 \(k[t]/(f)\)。取首一 \(f\) 后，零化子为 \((f)\)，不同的 \(f\) 给出不同同构类型。其 \(k\) 维数为 \(\deg f\)，因为 \(1,t,\ldots,t^{\deg f-1}\) 的陪集组成基。故一维情形正是 \(f=t-a\)。在实数域上，\(\mathbb R[t]/(t^2+1)\cong\mathbb C\) 为二维简单模，其中 \(t\) 作用为乘以 \(i\)。
\end{solution}

\begin{exercise}\label{b2:ex:13-squarefree-cyclic}
设 \(k\) 为域，\(V\ne0\) 为有限维 \(k\) 向量空间，\(T\in\operatorname{End}_k(V)\)，令 \(V_T\) 为由 \(t\cdot v=T(v)\) 定义的 \(k[t]\) 模。证明 \(V_T\) 半单当且仅当极小多项式 \(\mu_T(t)\) 不含重复不可约因子；\(T\) 在 \(k\) 上可对角化当且仅当 \(\mu_T\) 是互不相同的一次因子的乘积。分别分析实矩阵
\[
\begin{pmatrix}0&-1\\1&0\end{pmatrix},\qquad
\begin{pmatrix}0&1\\0&0\end{pmatrix}
\]
对应的模，说明半单性与在基域上可对角化的区别。
\end{exercise}
\begin{solution}
由\cref{b2:thm:rational-canonical-form,b2:prop:operator-characteristic-minimal}，存在首一非常数多项式 \(f_1\mid\cdots\mid f_s\)，使
\[
V_T\cong\bigoplus_{i=1}^s k[t]/(f_i),\qquad \mu_T=f_s.
\]
若 \(\mu_T\) 无重因子，每个 \(f_i\) 也无重因子，\cref{b2:prop:semisimple-cyclic-polynomial}使每个循环分量半单，故其直和半单。反过来，若 \(V_T\) 半单，则它的商 \(k[t]/(f_s)\) 半单，同一判据说明 \(f_s=\mu_T\) 无重因子。

若 \(T\) 可对角化，取由特征向量组成的基，每个基向量所在直线为简单 \(k[t]\) 模 \(k[t]/(t-\lambda)\)，零化全模的首一多项式中次数最小者正是全部不同特征值对应的 \(t-\lambda\) 之积。反过来，若 \(\mu_T\) 是不同一次因子的乘积，各 \(f_i\) 也如此；\secref{b2:sec:0403}的中国剩余分解把上述循环分量拆为一维模，合并这些一维模的基便给出特征向量基。

第一个实矩阵的极小多项式为 \(t^2+1\)：其平方为 \(-I\)，且它不是实标量矩阵。该多项式不可约，所以对应模简单、因而半单，但在 \(\mathbb R\) 上没有特征值，不能对角化。第二个矩阵非零而平方为零，极小多项式为 \(t^2\)，含重因子，因此对应模不半单。
\end{solution}

\begin{exercise}\label{b2:ex:13-graph-submodules}
设 \(R\) 为含幺环、\(S\) 为简单左 \(R\) 模，令 \(E=\operatorname{End}_R(S)\)。证明 \(S\oplus S\) 中与 \(S\oplus0\) 互补的子模恰为
\(K_f=\Bigl\{\,\bigl(f(s),s\bigr):s\in S\,\Bigr\}\)、\(f\in E\)。再分类 \(S\oplus S\) 的全部简单子模。
\end{exercise}
\begin{solution}
每个 \(K_f\) 是第二坐标投影的一个截面的像，且 \((x,y)=\bigl(x-f(y),0\bigr)+\bigl(f(y),y\bigr)\)，因此确为补模。反过来，对任一补模 \(K\)，第二投影 \(K\to S\) 单射，因为核包含于 \(K\cap(S\oplus0)\)；它满射来自直和覆盖。其逆形如 \(s\mapsto\bigl(f(s),s\bigr)\)，其中 \(f\) 线性，所以得到全部补模。

若简单子模 \(T\) 的第二投影为零，则 \(T\subseteq S\oplus0\)，由简单性 \(T=S\oplus0\)。若第二投影非零，Schur 引理使其成为同构，所以 \(T=K_f\)。这也说明一个固定直和分解不会列出所有简单子模。

更一般地，对任意正整数 \(n\)，\(\operatorname{Hom}_R(S,S^n)\cong E^n\) 是右 \(E\) 向量空间，非零同态均为单射；若非零同态 \(f,g:S\to S^n\) 的像相同，把它们视为到共同像的同构，则 \(e=g^{-1}f\in E^\times\) 且 \(f=g\circ e\)；反过来，这个等式也保证像相同。任一简单子模向某个坐标的投影非零，故由 Schur 引理同构于 \(S\)，所以全部简单子模都由这些同态得到。因此，\(S^n\) 的简单子模对应于这个右 \(E\) 向量空间的一维子空间，也就是右 \(E\) 空间的射影空间。
\end{solution}

\begin{exercise}\label{b2:ex:13-products-of-simple-modules}
令 \(p\) 遍历全部素数。证明整数模 \(\bigoplus_p\mathbb Z/p\mathbb Z\) 半单，但 \(\prod_p\mathbb Z/p\mathbb Z\) 不半单。若把全部因子都换成同一个 \(\mathbb Z/2\mathbb Z\)，直积是否半单？
\end{exercise}
\begin{solution}
直和按定义半单。每个简单整数模都为某个 \(\mathbb Z/p\mathbb Z\)，因此任意半单整数模中的元素只有有限支集，必有有限阶。不同素数的直积却含元素 \(x=(\bar1)_p\)，它不被任何非零整数 \(a\) 零化：取不整除 \(a\) 的素数 \(p\)，该坐标仍非零。因此直积不半单。若全部因子都是 \(\mathbb Z/2\mathbb Z\)，整个直积是 \(\mathbb F_2\) 向量空间，选取向量空间基后成为 \(\mathbb Z/2\mathbb Z\) 的直和，所以作为整数模半单。这个基不同于标准坐标向量族，后者只张成有限支集部分。同一个乘积作为其自身环的左正则模不半单，而作为整数模却半单；必须明确标量环，不能仅由“无限直积”判断半单性。
\end{solution}

\begin{exercise}\label{b2:ex:13-schur-converse}
令 \(M=\mathbb Q\oplus\mathbb Z\)，视为整数模。计算它的自同态环，并解释为什么从 \(\mathbb Q\) 到 \(\mathbb Z\) 的块必为零，而从 \(\mathbb Z\) 到 \(\mathbb Q\) 的块可以非零。求这个环中的全部幂等元。
\end{exercise}
\begin{solution}
由\cref{b2:prop:schur-hypothesis-counterexamples}，\(\operatorname{End}_{\mathbb Z}(\mathbb Q)=\mathbb Q\)；而整数模同态 \(\mathbb Z\to\mathbb Z\) 或 \(\mathbb Z\to\mathbb Q\) 都由 \(1\) 的像决定，分别得到 \(\mathbb Z\) 或 \(\mathbb Q\)。若 \(f:\mathbb Q\to\mathbb Z\)，则对每个正整数 \(n\)，有 \(f(1)=n f(1/n)\)，所以 \(f(1)\) 被所有正整数整除，只能为零。再由 \(b f(a/b)=a f(1)=0\) 和 \(\mathbb Z\) 无挠，得 \(f=0\)。因此
\[
\operatorname{End}_{\mathbb Z}(M)\cong
\left\{\begin{pmatrix}a&b\\0&d\end{pmatrix}:a,b\in\mathbb Q,\ d\in\mathbb Z\right\},
\]
其作用为 \((x,z)\mapsto(ax+bz,dz)\)，复合给出通常矩阵乘法。这里 \(\mathbb Q\) 与 \(\mathbb Z\) 并非一对简单模，不能借 Schur 引理把两个不同类型间的 Hom 都消去。

幂等条件为 \(a^2=a\)、\(d^2=d\)、\((a+d)b=b\)。前两式使 \(a,d\in\{0,1\}\)；若 \(a=d\)，则 \(b=0\)，得到零元和单位元。若 \((a,d)=(1,0)\) 或 \((0,1)\)，任意 \(b\in\mathbb Q\) 都可取，分别得到
\[
\begin{pmatrix}1&b\\0&0\end{pmatrix},\qquad
\begin{pmatrix}0&b\\0&1\end{pmatrix}.
\]
它们给出不同投影，表明即使两个固定分量间的一侧 Hom 为零，另一侧的非零 Hom 仍可改变补模。
\end{solution}

\begin{exercise}\label{b2:ex:13-schur-finite-dimension}
设 \(k\) 为域，将 \(M=k(t)\) 视为 \(k[t]\) 模，其中 \(t\) 按乘法作用。证明 \(\operatorname{End}_{k[t]}(M)\cong k(t)\)，但 \(M\) 不简单。比较把标量环扩大为 \(k(t)\) 后的简单性，并证明乘以 \(t\) 的自同态没有属于 \(k\) 的特征值。
\end{exercise}
\begin{solution}
对 \(f\in\operatorname{End}_{k[t]}(M)\)，令 \(c=f(1)\)。任取 \(a,b\in k[t]\)、\(b\ne0\)，线性性给出
\[
b f(a/b)=f(a)=ac.
\]
在 \(k(t)\) 中乘以 \(b\) 可逆，故 \(f(a/b)=(a/b)c\)。反过来，任意 \(c\in k(t)\) 的乘法映射都与 \(k[t]\) 作用相容，其加法和复合对应 \(k(t)\) 的加法与乘法，因此得到环同构。\(k[t]\) 是 \(M\) 的非零真 \(k[t]\) 子模，故 \(M\) 不简单；扩大为 \(k(t)\) 作用后，它成为域的一维向量空间，故简单。两种标量环给出相同的自同态环，却有不同的子模。

若乘以 \(t\) 有特征值 \(\lambda\in k\)，则存在 \(0\ne v\in k(t)\) 使 \((t-\lambda)v=0\)，与 \(k(t)\) 为域且 \(t-\lambda\ne0\) 矛盾。即使 \(k\) 代数闭，这个无限维算子也没有特征值；有限维 Schur 推论中的特征值步骤不能用于它。
\end{solution}

\begin{exercise}\label{b2:ex:13-endomorphism-blocks}
计算整数模 \(M=(\mathbb Z/2\mathbb Z)^3\oplus(\mathbb Z/3\mathbb Z)^2\) 的自同态环，并求其元素总数。
\end{exercise}
\begin{solution}
从二阶循环群到三阶循环群的同态像同时被二和三零化，所以只能为零，反方向也是如此。每个同类型分量中的同态由域中的一个标量决定，于是
\[
\operatorname{End}_{\mathbb Z}(M)\cong M_3(\mathbb F_2)\times M_2(\mathbb F_3).
\]
每个三阶矩阵有九个二值条目，每个二阶矩阵有四个三值条目，故元素总数为 \(2^9\cdot3^4\)。这里自同态环包含不可逆矩阵，因此它不是除环，符合 \(M\) 并非简单模的事实。
\end{solution}

\begin{exercise}\label{b2:ex:13-idempotent-summands}
设 \(k\) 为域，在 \(R=M_2(k)\) 中，令 \(I=RE_{11}\) 为第二列全零的左理想。求全部以 \(I\) 为像且在 \(I\) 上为恒等的左模投影 \(p:R\to I\)，写出各自的幂等元及补模。一个直和因子是否唯一决定它在正则模中的投影？
\end{exercise}
\begin{solution}
由\cref{b2:prop:regular-hom-endomorphism}，投影必为右乘 \(e=p(1)\in I\)，写 \(e=\begin{psmallmatrix}a&0\\b&0\end{psmallmatrix}\)。要求 \(p(E_{11})=E_{11}\) 给出 \(a=1\)；对任意 \(b\in k\)，所得
\[
e_b=\begin{pmatrix}1&0\\b&0\end{pmatrix}
\]
都幂等，并且右乘它在 \(I\) 上为恒等。若以 \((v,w)\) 表示一个矩阵的两列，则 \(p_b(v,w)=(v+bw,0)\)，故核为所有 \((-bw,w)\)、\(w\in k^2\) 组成的左理想。\cref{b2:prop:idempotent-module-decomposition}给出 \(R=I\oplus\ker p_b\)。不同 \(b\) 给出不同投影及补模，所以只指定直和因子 \(I\) 还不能唯一指定投影。
\end{solution}

\begin{exercise}\label{b2:ex:13-annihilator-vector-module}
设 \(D\) 为除环，\(n\ge1\)，令 \(R=M_n(D)\)、\(S=D^n\) 为标准幺性左列模。求 \(\operatorname{Ann}_R(S)\) 和标准列向量 \(e_1\) 的零化左理想，指出二者何时相同、何时不同。
\end{exercise}
\begin{solution}
若矩阵零化所有标准列，它的每列都为零，因此 \(\operatorname{Ann}_R(S)=0\)。而 \(Ae_1=0\) 当且仅当 \(A\) 的第一列为零，所以单个向量的零化子是这些矩阵组成的左理想。\(n=1\) 时这个理想也是零，两个零化子相同；对 \(n>1\)，它非零且不是双边理想，例如 \(E_{12}\) 的第一列为零，但 \(E_{12}E_{21}=E_{11}\) 的第一列非零。零化整个模得到双边理想，零化一个元素一般只得到左理想。这里 \(S\) 又是简单列模，且 \(\operatorname{Ann}_R(S)=0\)，所以 \(S\) 是忠实简单模；\chapref{b2:ch:14}将以忠实简单模定义本原环。
\end{solution}

\begin{exercise}\label{b2:ex:13-center-and-idempotents}
设 \(k\) 为含 \(q\) 个元素的有限域，\(R=M_2(k)\times k\)。分类 \(R\) 中幂等元的共轭类型，分别计算全部幂等元和中心幂等元的个数。这里共轭指由 \(R\) 中单位元实施的 \(e\mapsto ueu^{-1}\)。
\end{exercise}
\begin{solution}
若 \(P\in M_2(k)\) 幂等，线性投影给出 \(k^2=\operatorname{im}P\oplus\ker P\)。选择两部分的基，便使 \(P\) 共轭于 \(\operatorname{diag}(I_r,0)\)，其中 \(r=0,1,2\)，而秩在共轭下不变，所以恰有三种共轭类型。\(k\) 的幂等元只有零和一，且共轭不改变它们。因此 \(R\) 中的幂等元共有六种共轭类型。

秩一的投影由像直线 \(L\) 与核直线 \(K\ne L\) 唯一确定。\(k^2\) 有 \((q^2-1)/(q-1)=q+1\) 条直线，选定 \(L\) 后有 \(q\) 种 \(K\)，故秩一幂等矩阵共 \(q(q+1)\) 个。加上零矩阵和单位矩阵，再为 \(k\) 因子独立选择零或一，得到全部幂等元的个数
\[
2\bigl(2+q(q+1)\bigr).
\]
由\cref{b2:prop:matrix-centers-central-idempotents}，中心幂等元却只有 \((0,0),(I_2,0),(0,1),(I_2,1)\) 四个。秩一投影不在中心，不能把它切出的左模分量误当成乘积环的双边矩阵块。
\end{solution}

\begin{exercise}\label{b2:ex:13-concrete-finite-product}
设 \(R=M_2(\mathbb F_3)\times M_3(\mathbb F_2)\)。求全部简单左模的同构类型、各自元素个数、它们在左正则模中的重数，以及 \(R\) 的元素总数。
\end{exercise}
\begin{solution}
两种简单模分别为 \(S=\mathbb F_3^2\) 和 \(T=\mathbb F_2^3\)，由对应因子作用，另一因子作用为零。它们分别有九个、八个元素；通过中心幂等元的不同作用可知不可能同构。按列分解，\({}_RR\cong S^2\oplus T^3\)，所以重数为二和三，正则模长度为五。环本身的元素总数为 \(3^4\cdot2^9\)。两种模来自不同特征，这里不为它们虚构共同的基域。
\end{solution}

\begin{exercise}\label{b2:ex:13-dimension-nine}
设 \(k\) 为代数闭域。列出所有九维结合含幺半单 \(k\) 代数的同构类型，并为每个类型列出简单模的维数。
\end{exercise}
\begin{solution}
由\cref{b2:prop:finite-dimensional-semisimple-algebra}，维数是矩阵大小的平方之和。可能的正平方为 \(1,4,9\)，故全部类型为
\[
k^9,\qquad M_2(k)\times k^5,\qquad
M_2(k)\times M_2(k)\times k,\qquad M_3(k).
\]
相应简单模维数依次为九个一；一个二和五个一；两个二和一个一；一个三。Artin–Wedderburn 参数唯一性说明这些类型彼此不同，也说明没有遗漏。
\end{solution}

\begin{exercise}\label{b2:ex:13-commutative-semisimple}
证明交换含幺半单环恰为有限个域的乘积，包括零环的空乘积情形。指出为什么这里不能把“有限个”改为“任意多个”。
\end{exercise}
\begin{solution}
在 Artin–Wedderburn 分解中，若某个 \(n_i>1\)，则 \(E_{11}E_{12}=E_{12}\) 而 \(E_{12}E_{11}=0\)，不交换，所以所有 \(n_i=1\)。各 \(D_i\) 作为该环的商也交换，因此是域。反过来，有限个域的乘积由结构定理半单且显然交换。无限乘积 \(\prod_{j\ge1}k\) 中“前 \(r\) 个坐标为零”的理想给出无限严格降链，而半单环必须 Artin，所以不满足结论。
\end{solution}

\begin{exercise}\label{b2:ex:13-cyclic-three}
分别在 \(\mathbb Q\)、\(\mathbb C\) 及特征三的域 \(k\) 上分析循环群 \(C_3\) 的群代数。前两种情形写出半单分解；后一种情形证明不半单。不要调用一般 Maschke 定理。
\end{exercise}
\begin{solution}
群代数为 \(k[t]/(t^3-1)\)。在 \(\mathbb Q\) 上，\(t^3-1=(t-1)(t^2+t+1)\)，第二个因子没有有理根所以不可约，且两因子互素。\secref{b2:sec:0403}的中国剩余分解给出 \(\mathbb Q[C_3]\cong\mathbb Q\times\mathbb Q(\zeta_3)\)，其中 \(\zeta_3\) 为本原三次单位根，满足 \(\zeta_3^2+\zeta_3+1=0\)，故第二因子就是 \(\mathbb Q[t]/(t^2+t+1)\)。在 \(\mathbb C\) 上，三个单位根互异，故 \(\mathbb C[C_3]\cong\mathbb C^3\)。

特征三时 \(t^3-1=(t-1)^3\)，写 \(\varepsilon=t-1\) 得 \(R=k[\varepsilon]/(\varepsilon^3)\)。若 \((\varepsilon^2)\) 有模投影 \(p:R\to(\varepsilon^2)\)，则 \(p(\varepsilon^2)=\varepsilon^2p(1)=0\)，但投影在该非零理想上应为恒等，矛盾。故此环不半单。
\end{solution}

\begin{exercise}\label{b2:ex:13-finite-order-p-four}
设 \(p\) 为素数。分类所有恰有 \(p^4\) 个元素的含幺半单环，并指出哪些非交换。
\end{exercise}
\begin{solution}
由\cref{b2:cor:finite-semisimple-rings}，所有因子形如 \(M_n(\mathbb F_{p^f})\)，它对元素总数的指数贡献为 \(fn^2\)，各贡献之和为四。若 \(n>1\)，只能 \(n=2,f=1\)，且此因子已占全部指数，得到 \(M_2(\mathbb F_p)\)。其余情况全为域，按四的正整数分拆得到
\[
\mathbb F_{p^4},\quad \mathbb F_{p^3}\times\mathbb F_p,\quad
\mathbb F_{p^2}\times\mathbb F_{p^2},\quad
\mathbb F_{p^2}\times\mathbb F_p\times\mathbb F_p,\quad
\mathbb F_p^4.
\]
只有矩阵环这一种非交换。第一卷定理17.1.1与定理17.1.2保证各素数幂大小的有限域存在，且由元素数决定同构类型；再加上结构定理的唯一性，得到列表无重叠且穷尽。
\end{solution}

\begin{exercise}\label{b2:ex:13-smallest-noncommutative-simple}
分类所有有 \(p^8\) 个元素的含幺单环，其中 \(p\) 为素数。再求非交换有限单环可能具有的最小元素数。
\end{exercise}
\begin{solution}
有限环当然左 Artin，因此单环由\cref{b2:prop:simple-artinian}写为 \(M_n(D)\)。\(D\) 有限，由 Wedderburn 小定理为 \(\mathbb F_{p^f}\)，且 \(fn^2=8\)。于是只有 \(n=1,f=8\) 或 \(n=2,f=2\)，即 \(\mathbb F_{p^8}\) 与 \(M_2(\mathbb F_{p^2})\)。一般非交换有限单环必须 \(n\ge2\)，所以元素数至少为 \(2^{2^2}=16\)，且 \(M_2(\mathbb F_2)\) 达到这个值。
\end{solution}

\begin{exercise}\label{b2:ex:13-recover-multiplicities}
设 \(k\) 为域，\(R=M_2(k)\times k\)，\(M\) 为有限维幺性左 \(R\) 模，且 \(\dim_kM=8\)、\(\dim_k\operatorname{End}_R(M)=13\)。求 \(M\) 的简单分解。
\end{exercise}
\begin{solution}
两种简单模为 \(S=k^2\)、\(T=k\)，自同态环均为 \(k\)，不同类型间 Hom 为零。由\cref{b2:thm:artin-wedderburn}，\(R\) 是半单环，\cref{b2:thm:semisimple-ring-characterizations}于是保证每个左 \(R\) 模半单，包括题设的 \(M\)。因此 \(M\cong S^a\oplus T^b\)，其中 \(a,b\ge0\)。同类分量间的同态给出矩阵块，不同类型间的同态为零，故 \(\operatorname{End}_R(M)\cong M_a(k)\times M_b(k)\)，其维数为 \(a^2+b^2\)。题设两项维数因此给出
\(2a+b=8\)、\(a^2+b^2=13\)。第二式迫使 \((a,b)=(2,3)\) 或 \((3,2)\)，第一式选出 \((3,2)\)。因此 \(M\cong S^3\oplus T^2\)。这两项数值共同恢复重数，单用总维数则不能做到。
\end{solution}

\begin{exercise}\label{b2:ex:13-dual-numbers-semisimple-modules}
设 \(k\) 为域，\(R=k[\varepsilon]/(\varepsilon^2)\)。证明一个幺性左 \(R\) 模半单，当且仅当 \(\varepsilon\) 在其上作用为零。把结论翻译成满足 \(T^2=0\) 的线性算子的陈述。
\end{exercise}
\begin{solution}
若 \(S\) 简单，\(\varepsilon S\) 是子模，因为环交换。若它非零，则等于 \(S\)，于是 \(\varepsilon^2S=\varepsilon S=S\)，与 \(\varepsilon^2=0\) 矛盾。所以 \(\varepsilon S=0\)，每个半单模的全部简单分量都被 \(\varepsilon\) 零化，整个模也如此。反过来，若 \(\varepsilon M=0\)，作用经过 \(R/(\varepsilon)=k\)，选取 \(k\) 基后每个一维子空间都是简单 \(R\) 子模，故半单。若用 \(T\) 表示 \(\varepsilon\) 的作用，结论就是对应模半单当且仅当 \(T=0\)，而不仅仅是 \(T^2=0\)。
\end{solution}

\begin{exercise}\label{b2:ex:13-infinite-product-semiprimitivity}
令 \(R=\prod_{i\ge1}k_i\)，其中各 \(k_i\) 为域。对每个 \(i\)，令 \(S_i=k_i\) 通过第 \(i\) 个坐标接受 \(R\) 作用。证明 \(R\) 满足半本原性的零化子条件，却不是半单环。
\end{exercise}
\begin{solution}
各 \(S_i\) 简单，因为其非零向量在对应域作用下生成整个 \(k_i\)。\(\operatorname{Ann}_R(S_i)\) 正是第 \(i\) 个坐标为零的理想，这些零化子的交为零。因此所有简单模零化子的交作为其子集也为零，\(R\) 半本原。另一方面，依次要求前一个、前两个、前三个坐标为零，得到无限严格下降的左理想链。由\cref{b2:thm:semisimple-ring-characterizations}，半单环必左 Artin，故此 \(R\) 不半单。
\end{solution}

\begin{exercise}\label{b2:ex:13-projective-factors}
设 \(R\) 为含幺环，幺性左 \(R\) 模 \(M\) 有有限合成列，且该合成列中每个简单商因子都是投射模。证明 \(M\) 半单。反过来，半单模的简单因子是否必为原环上的投射模？
\end{exercise}
\begin{solution}
对合成列长度归纳。长度零时为零模。设末端为 \(0\to N\to M\xrightarrow{q}S\to0\)，其中 \(S\) 为投射简单模。将 \(\operatorname{id}_S\) 沿满射 \(q\) 提升，得到截面，故 \(M\cong N\oplus S\)。\(N\) 的较短合成列仍有投射简单商因子，归纳得它半单，所以 \(M\) 半单。逆命题不成立：取素数 \(p\)，\(\mathbb Z/p\mathbb Z\) 本身简单且半单，但整数模满射 \(\mathbb Z\to\mathbb Z/p\mathbb Z\) 不分裂，因此它不是投射模。对有限长度模，沿一条合成列逐层分裂，就能把全部简单商提升为实际直和分量；反过来，半单模的子模都有补模，合成列中的各次扩张都分裂。简单商因子投射是保证这些扩张分裂的充分条件，却不是半单性所必需的条件。
\end{solution}
